\section*{अध्याय \ref{ch:g7:identifying-substances} --- पदार्थों की पहचान}

\begin{solution}{exo:g7:identifying-substances:1}
चूने के पानी का परीक्षण: गैस को स्वच्छ चूने के पानी में से बुलबुलों के रूप
में गुज़ारा जाता है; अगर गैस कार्बन डाइऑक्साइड है, तो वह दूधिया हो जाता है।
\end{solution}

\begin{solution}{exo:g7:identifying-substances:2}
चूर्ण निर्जल कॉपर सल्फ़ेट है, और द्रव में पानी है।
\end{solution}

\begin{solution}{exo:g7:identifying-substances:3}
ऑक्सीजन।
\end{solution}

\begin{solution}{exo:g7:identifying-substances:4}
पेंसिल \omterm{def:g6:solutions-and-solubility:solute}{विलायक} में नहीं \omterm{def:g2:mixing-and-dissolving:dissolve}{घुलती}। स्याही की रेखा ख़ुद ऊपर ले जाई जाती और अलग
हो जाती, और \omterm{def:g7:identifying-substances:chromatography}{वर्णलेख} को बिगाड़ देती।
\end{solution}

\begin{solution}{exo:g7:identifying-substances:5}
अगर धब्बा \omterm{def:g6:solutions-and-solubility:solute}{विलायक} के नीचे होता, तो रंजक कागज़ पर चढ़ने के बजाय बीकर के
\omterm{def:g6:solutions-and-solubility:solute}{विलायक} में घुल जाते।
\end{solution}

\begin{solution}{exo:g7:identifying-substances:6}
नहीं, वह दो रंजकों का \omterm{def:g6:pure-substances-and-mixtures:mixture}{मिश्रण} है: शायद संदर्भों वाला पीला रंजक और नीला
रंजक।
\end{solution}

\begin{solution}{exo:g7:identifying-substances:7}
गैस को थोड़ी मात्रा में एक परखनली में इकट्ठा करो और उसके मुँह पर \omterm{def:g4:what-a-fire-needs:burning}{जलती}
खपच्ची पकड़ो: तीखी `पॉप' की आवाज़ हाइड्रोजन दिखाती है।
\end{solution}

\begin{solution}{exo:g7:identifying-substances:8}
बाहर छोड़ी गई साँस में ताज़ी \omterm{def:g4:air-a-mixture-of-gases:air}{हवा} से कहीं ज़्यादा कार्बन डाइऑक्साइड होती है।
\end{solution}

\begin{solution}{exo:g7:identifying-substances:9}
परीक्षण दिखाता है कि पानी मौजूद है, यह नहीं कि वह अकेला है: नमकीन पानी,
रस या पानी वाला कोई भी \omterm{def:g6:pure-substances-and-mixtures:mixture}{मिश्रण} चूर्ण को नीला कर देगा।
\end{solution}

\begin{solution}{exo:g7:identifying-substances:10}
तीन रंजक (तीन धब्बे)। पीला रंजक सबसे ऊपर चढ़ा।
\end{solution}

\begin{solution}{exo:g7:identifying-substances:11}
उदाहरण के लिए: हर जार में सुलगती खपच्ची डालो --- जो उसे फिर जला दे, उसमें
ऑक्सीजन है। बाकी दोनों के मुँह पर \omterm{def:g4:what-a-fire-needs:burning}{जलती} खपच्ची पकड़ो --- तीखी `पॉप'
हाइड्रोजन दिखाती है। आख़िरी जार में कार्बन डाइऑक्साइड है (चूने के पानी से इसकी
पुष्टि की जा सकती है)। दो या तीन परीक्षण काफ़ी हैं।
\end{solution}

\begin{solution}{exo:g7:identifying-substances:12}
हो सकता है गैस में कार्बन डाइऑक्साइड हो ही नहीं, या उस समय में दिखने लायक
बहुत कम हो। बाहर छोड़ी गई साँस से, जिसमें कार्बन डाइऑक्साइड होना निश्चित है,
किया गया नियंत्रण परीक्षण जाँचता है कि चूने का पानी ठीक काम कर रहा है।
\end{solution}

\begin{solution}{pb:g7:identifying-substances:1}
\textbf{1.} ताकि चारों के लिए परिस्थितियाँ (कागज़, \omterm{def:g6:solutions-and-solubility:solute}{विलायक}, समय) एक जैसी
हों: तभी धब्बों की ऊँचाइयों की तुलना की जा सकती है।

\textbf{2.} \omterm{def:g6:pure-substances-and-mixtures:mixture}{मिश्रण}: उसकी स्याही तीन अलग धब्बे देती है, यानी तीन रंजक।

\textbf{3.} दो रंजक। नहीं: पर्चे की स्याही में तीन रंजक हैं, और पेन 1 में
न लाल रंजक है, न पीला।

\textbf{4.} उसके तीनों धब्बे पर्चे के धब्बों की ऊँचाइयों पर नहीं हैं: वे
तीन अलग रंजक हैं।

\textbf{5.} पेन 2: उसके तीनों धब्बे ठीक पर्चे के तीनों धब्बों की
ऊँचाइयों पर हैं।

\textbf{6.} ऑक्सीजन।

\textbf{7.} हाइड्रोजन।

\textbf{8.} कार्बन डाइऑक्साइड।

\textbf{9.} रिक्त परीक्षण (एक नियंत्रण)। यह दिखाता है कि उस समय में चूने का
पानी न अपने आप दूधिया होता है, न साधारण \omterm{def:g4:air-a-mixture-of-gases:air}{हवा} से: इसलिए गैस Z के साथ आया
दूधियापन सचमुच कार्बन डाइऑक्साइड से है।

\textbf{10.} ``पानी --- निर्जल कॉपर सल्फ़ेट --- सफ़ेद से नीला'':
द्रव में पानी है।

\textbf{11.} नहीं। धनात्मक परीक्षण दिखाता है कि पानी मौजूद है, यह नहीं कि और
कुछ नहीं है: द्रव कोई विलयन या पानी वाला कोई दूसरा \omterm{def:g6:pure-substances-and-mixtures:mixture}{मिश्रण}
हो सकता है।

\textbf{12.} पर्चे की स्याही में 3 रंजक हैं, और पर्चा पेन 2 ने
लिखा।
\end{solution}
